\chapter{Hip Pain}

\section*{Definition}
Discomfort is localized to the hip joint or surrounding structures (groin, lateral thigh, buttock), arising from intra-articular pathology, periarticular soft tissues, or referred sources.

\section*{What Happens in Your Body}
Hip pain localizes differently based on origin: groin pain typically indicates intra-articular pathology (osteoarthritis, labral tear, AVN); lateral thigh pain suggests trochanteric bursitis; buttock pain may arise from the sacroiliac joint, lumbar spine, or posterior hip structures.

\section*{Causes}
\begin{itemize}
  \item Hip osteoarthritis (degenerative joint disease)
  \item Trochanteric bursitis
  \item Hip labral tear
  \item Avascular necrosis (AVN) of femoral head
  \item Femoroacetabular impingement (FAI)
  \item Hip fracture (especially in elderly/osteoporotic)
  \item Iliopsoas tendinopathy or bursitis
  \item Referred low back pain or sacroiliac dysfunction
  \item Osteitis pubis
  \item Septic arthritis (urgent, with fever and severe pain)
\end{itemize}

\section*{When to See a Doctor}
See your doctor if:
\begin{itemize}
  \item Symptoms are severe or getting worse over time
  \item Severe hip pain after fall or injury, inability to bear weight, fever with hip pain, or groin pain with fever
  \item You have other concerning symptoms such as fever, weight loss, or bleeding
\end{itemize}

\section*{Related Symptoms}
\begin{itemize}
  \item Groin pain
  \item Limping
  \item Stiffness
  \item Lower back pain
  \item Knee pain
\end{itemize}
\textbf{Important:} If this symptom persists, your doctor will check for serious underlying causes.

\section*{Self-Care / Home Management}
\begin{itemize}
  \item Avoid activities that aggravate hip pain such as prolonged walking on uneven terrain, running, or deep squatting; switch to low-impact exercise like swimming or cycling.
  \item Apply ice packs to the lateral hip or groin for 15-20 minutes after activity if pain is acute; use a heating pad for chronic stiffness.
  \item Use a walking aid (cane held in the opposite hand) temporarily if weight-bearing is painful to reduce joint load.
  \item Maintain a healthy body weight to decrease mechanical stress on the hip joint.
\end{itemize}

\section*{How Your Doctor Will Evaluate You}
\begin{itemize}
  \item Assessment of gait, hip range of motion, strength, and specific provocative tests (FABER, FADIR, log roll) to localise the source of pain.
  \item Plain radiographs (weight-bearing pelvis and lateral hip) are first-line imaging to evaluate for osteoarthritis, AVN, or fracture.
  \item MRI with intra-articular contrast is indicated when labral tear, impingement, or early AVN is suspected and X-rays are normal.
  \item Referral to an orthopaedic surgeon or sports medicine specialist is made for persistent symptoms, mechanical symptoms (catching, locking), or suspected surgical pathology.
\end{itemize}

\section*{Treatment Options}
\begin{itemize}
  \item Physical therapy targeting hip abductor and core strengthening, along with activity modification, is first-line treatment for most non-arthritic hip pain.
  \item NSAIDs or acetaminophen provide symptomatic relief; intra-articular corticosteroid injections can reduce inflammation in osteoarthritis or bursitis.
  \item Surgical options include hip arthroscopy (labral repair, FAI decompression), osteotomy, and total hip arthroplasty for advanced osteoarthritis or AVN.
  \item Weight-bearing precautions and orthopaedic follow-up are essential after hip fracture or arthroplasty to prevent dislocation and ensure proper rehabilitation.
  \item Medications, lifestyle modifications, and physical therapies may be prescribed as appropriate.
  \item Follow-up ensures treatment effectiveness and allows timely adjustments to the care plan.
\end{itemize}

\section*{References}

\begin{thebibliography}{99}
\bibitem{harrison-hip} Longo DL, et al. \emph{Harrison's Principles of Internal Medicine}. 21st ed. McGraw-Hill; 2022. Chapter 391: Approach to the Patient with Hip Pain.
\bibitem{uptodate-hip} Lane NE, et al. Evaluation of hip pain in adults. In: UpToDate, Post TW (Ed), UpToDate, Waltham, MA; 2025.
\bibitem{aaos-hip} American Academy of Orthopaedic Surgeons. \emph{AAOS Clinical Practice Guideline: Management of Osteoarthritis of the Hip}. AAOS; 2023.
\bibitem{tanzer} Tanzer M, et al. Hip osteoarthritis: a clinical review. \emph{JAMA}. 2023;329(10):832-843.
\bibitem{canale} Azar FM, et al. \emph{Campbell's Operative Orthopaedics}. 14th ed. Elsevier; 2023. Chapter 5: Hip Disorders.
\bibitem{klauser} Klauser AS, et al. Imaging of hip pain: a systematic review. \emph{Radiology}. 2023;306(3):e221494.
\end{thebibliography}

\clearpage